\section{Failure recovery and implementation discipline}\label{sec:C12:recovery}
Context transformations should be recoverable operations. Persist the source view and its edit manifest before replacing the active context. Commit the replacement only after structural validation, and retain a rollback reference until the affected task stage has been verified. This does not require keeping every cache allocation forever; it requires retaining enough evidence and metadata to reconstruct the decision input. Storage retention and privacy constraints still apply.

Use separate error classes for failed summarization, invalid serialization, missing raw evidence, retrieval failure, and an edited view that violates a pinned obligation. A generic ``context error'' obscures whether the model, controller, or storage layer is responsible. Recovery can then be targeted: restore the prior view, fetch an omitted record, regenerate an abstraction, or stop before a consequential tool effect. Blindly retrying the same damaged summary can amplify the failure.

Concurrent editing requires ownership. If a worker updates a plan while a controller compacts an older version, a last-write-wins replacement can erase legitimate progress. An optimistic commit should compare the base view version with the current version and reject or merge a stale edit. Semantic merges deserve special treatment: two summaries can each be locally valid while their union contains contradictory commitments. A storage-level conflict resolution rule cannot determine which commitment the user currently authorizes.

Human-readable state remains useful for audit, but readability should not replace machine-checkable boundaries. Identifiers, timestamps, tool-effect receipts, authority, and completion predicates can be represented in typed fields. Generated prose can explain those fields without becoming their sole source of truth. The boundary between recorded fact and model inference should survive every compaction. This reduces the chance that repeated summaries gradually promote a tentative hypothesis into an asserted observation.

The implementation path is incremental. First instrument the existing context assembly. Then introduce explicit source references and pinned obligations. Next add reversible offload with deterministic triggers. Only after those controls are measurable should a learned editing policy be compared under matched budgets. This sequence is a systems deduction about observability and attribution, not a claim that one particular product follows it or that it guarantees success.

The strongest outcome of context engineering is not a permanently small prompt. It is a working view whose contents, omissions, cost, and authority are deliberate and inspectable. Such a view supports the policy's current decisions while preserving a path back to evidence when assumptions change. Durable sessions in Chapter~\ref{ch:C13} provide the substrate for that path; persistent memory in Chapter~\ref{ch:C14} determines which evidence and abstractions remain useful across tasks.
